% STARTHERE

La struttura seguente ricalca la guideline di reporting STARD (i
numeri tra parentesi graffe, che si consiglia di lasciare, indicano
gli item della stessa); ai fini della compilazione delle parti di
propria competenza si consiglia di fare riferimento allo
\href{https://www.ncbi.nlm.nih.gov/pmc/articles/pmid/28137831/}{STARD
  explanation and elaboration}

\section{Title \{1\}}

{\em Identification as a study of diagnostic accuracy using at least one
measure of accuracy (such as sensitivity, specificity, predictive
values, or AUC)}

\section{Abstract \{2\}}

{\em Structured summary of study design, methods, results, and conclusions
(for specific guidance, see STARD for Abstracts)}

\section{Introduction}

{\em Scientific and clinical background, including the intended use and
clinical role of the index test}

\subsection{Background \{3\}}

{\em Scientific and clinical background, including the intended use and
clinical role of the index test}

\subsection{Objectives and hypotheses \{4\}}

{\em Study objectives and hypotheses}

\section{Methods}

\subsection{Study design \{5\}}

{\em Whether data collection was planned before the index test and
reference standard were performed (prospective study) or after
(retrospective study)}

\subsection{Partecipants}

\subsubsection{Eligibility criteria \{6\}}

{\em Eligibility criteria}

\subsubsection{Potential partecipants identification \{7\}}

{\em On what basis potentially eligible participants were identified (such
as symptoms, results from previous tests, inclusion in registry)}

\subsubsection{Setting, location and dates of partecipants
  identification \{8\}}

{\em Where and when potentially eligible participants were identified
  (setting, location and dates)}

\subsubsection{Sampling \{9\}}

{\em Whether participants formed a consecutive, random or convenience series}

\subsection{Test methods}

\subsubsection{Index test \{10a\}}

{\em Index test, in sufficient detail to allow replication}

\subsubsection{Reference standard \{10b\}}

{\em Reference standard, in sufficient detail to allow replication}

\subsubsection{Rationale for the reference standard \{11\}}

{\em Rationale for choosing the reference standard (if alternatives exist)}

\subsubsection{Definition/rationale for cutoffs/categories of the
  index test \{12a\}}

{\em Definition of and rationale for test positivity cut-offs or
  result categories of the index test, distinguishing pre-specified
  from exploratory}

\subsubsection{Definition/rationale for cutoffs/categories of the
  reference standard \{12b\}}

{\em Definition of and rationale for test positivity cut-offs or
  result categories of the reference standard, distinguishing
  pre-specified from exploratory}

\subsubsection{Blindness of index test performers/readers \{13a\}}

{\em Whether clinical information and reference standard results were
available to the performers/readers of the index test}

\subsubsection{Blindness of reference standard assessors \{13b\}}

{\em Whether clinical information and index test results were available
to the assessors of the reference standard}

\subsection{Analysis}

\subsubsection{Methods for diagnostic accuracy estimation \{14\}}

{\em Methods for estimating or comparing measures of diagnostic accuracy}

\subsubsection{Indeterminate results handling \{15\}}

{\em How indeterminate index test or reference standard results were handled}

\subsubsection{Missing data handling \{16\}}

{\em How missing data on the index test and reference standard were handled}

\subsubsection{Analyses of variability \{17\}}

{\em Any analyses of variability in diagnostic accuracy,
  distinguishing pre-specified from exploratory}

\subsubsection{Sample size \{18\}}

{\em Intended sample size and how it was determined}


\section{Results}


\subsection{Partecipants}

\subsubsection{Flow of participants \{19\}}

{\em Flow of participants, using a diagram}

\subsubsection{Participants characteristics \{20\}}

{\em Baseline demographic and clinical characteristics of participants}

\subsubsection{Severity of disease in patients with target conditions \{21a\}}

{\em Distribution of severity of disease in those with the target condition}

\subsubsection{Alternative diagnoses in those without the target
  condition \{21b\}}

{\em Distribution of alternative diagnoses in those without the target
  condition}

\subsubsection{Time/interventions between index test and reference
  standard \{22\}}

{\em Time interval and any clinical interventions between index test
  and reference standard}

\subsection{Test results}

\subsubsection{Cross tabulation \{23\}}

{\em Cross tabulation of the index test results (or their distribution) by
the results of the reference standard}

\subsubsection{Diagnostic accuracy estimates \{24\}}

{\em Estimates of diagnostic accuracy and their precision (such as 95\%
confidence intervals)}

\subsubsection{Adverse events \{25\}}

{\em Any adverse events from performing the index test or the
  reference standard}

\section{Discussion}

\subsection{Study limitations \{26\}}

{\em Study limitations, including sources of potential bias,
  statistical uncertainty, and generalisability}

\subsection{Implication for practice \{27\}}

{\em Implications for practice, including the intended use and
  clinical role of the index test}

\section{Other information}

\subsection{Protocol registration \{28\}}

{\em Registration number and name of registry}

\subsection{Protocol availability \{29\}}

{\em Where the full study protocol can be accessed}

\subsection{Funding \{30\}}

{\em Sources of funding and other support; role of funders}

